\subsection{Strategy comparison across data scales}

每个数据规模固定由dynamics validation选择出的同一个World Model，再比较CEM、iCEM、MPPI与MB-PPO。所有方法使用相同development seeds 42--46、相同1000步episode和相同reward definition。因此，\cref{fig:strategy_matrix,tab:strategy_matrix}中的差异主要来自strategy与固定World Model之间的交互，而不是更换模型或seed。

\begin{figure}[H]
  \centering
  \includegraphics[width=\textwidth]{fig5_strategy_comparison.pdf}
  \caption{\textbf{Development strategy comparison across data scales.} \textbf{a--c}, M100、M1000和M10000。空心点为seeds 42--46，实心点及误差条为mean $\pm$ population SD。M100 iCEM明显超出主轴范围，因此在inset中保留全部五个seed。CEM在M100最高；MPPI在M1000和M10000最高。}
  \label{fig:strategy_matrix}
\end{figure}


\begin{table}[H]
\centering
\caption{Development strategy matrix (seeds 42--46). 数值为 episode cumulative reward 的 mean $\pm$ population SD，越高越好。}
\label{tab:strategy_matrix}
\scriptsize
\resizebox{\textwidth}{!}{%
\begin{tabular}{lrrrr}
\toprule
Dataset & CEM & iCEM & MPPI & MB-PPO \\
\midrule
M100 & -274,238 $\pm$ 588 & -2,646,910 $\pm$ 4,581 & -348,211 $\pm$ 5,687 & -281,009 $\pm$ 131 \\
M1000 & -267,826 $\pm$ 407 & -334,597 $\pm$ 4,503 & -221,072 $\pm$ 1,475 & -285,326 $\pm$ 158 \\
M10000 & -269,218 $\pm$ 411 & -321,264 $\pm$ 36,310 & -223,774 $\pm$ 578 & -285,217 $\pm$ 283 \\
\bottomrule
\end{tabular}}
\end{table}


M100上的iCEM和MPPI均明显落后，尤其iCEM出现一致的极低真实回报；相同算法通过了已知dynamics和cost上的deterministic optimization test，因此这里更合理的解释是弱World Model与optimizer之间发生model exploitation，而不是直接断言算法不适合Industrial Benchmark。M1000和M10000中MPPI的提升说明World Model质量是重要瓶颈。

\subsection{World Model $\times$ Strategy controlled ablation}

M1000进一步完成两组控制变量实验：第一组固定MPPI配置，仅改变五种World Model；第二组固定MB-PPO训练协议，仅改变World Model。该实验研究dynamics prediction ranking与downstream control ranking是否完全一致，不用于反向修改validation-only模型选择。

\begin{figure}[H]
  \centering
  \includegraphics[width=\textwidth]{fig6_world_model_strategy_cross.pdf}
  \caption{\textbf{World Model architecture affects planning and learned policies differently.} \textbf{a}, 固定MPPI时五种World Model的development return。\textbf{b}, 固定MB-PPO protocol时的结果。每组$n=5$ seeds，点和误差条分别表示individual seed与mean $\pm$ population SD。MPPI在Transformer-4L上最高，而MB-PPO在LSTM上最高；dynamics validation最优仍为Transformer-2L。}
  \label{fig:cross}
\end{figure}

\subsection{Final simulator evaluation}

所有模型、strategy和设置冻结后，seeds 100--109才首次用于最终评价。最终系统按development结果确定：M100采用GRU+CEM，M1000与M10000采用Transformer-2L+MPPI。\Cref{fig:final,tab:final}同时展示原始行为数据参考值、BC与三个最终系统。

\begin{figure}[H]
  \centering
  \includegraphics[width=\textwidth]{fig7_final_simulator.pdf}
  \caption{\textbf{Final NeoRL simulator results on untouched seeds.} Original Behavior以灰色box表示NeoRL数据集1000条已有轨迹的累计奖励分布，属于unpaired dataset-trajectory reference。其余方法均部署到NeoRL simulator；空心点为seeds 100--109，实心点及误差条为mean $\pm$ population SD。三个最终系统在十个seed上均优于BC。}
  \label{fig:final}
\end{figure}


\begin{table}[H]
\centering
\caption{Final NeoRL evaluation. Online systems use untouched seeds 100--109. Original Behavior来自数据集已有轨迹，不是同seed在线部署。}
\label{tab:final}
\scriptsize
\resizebox{\textwidth}{!}{%
\begin{tabular}{llrrr}
\toprule
Method/System & Evaluation source & Episode return & $\Delta$ vs BC & Win rate \\
\midrule
原始行为数据参考值（Original Behavior） & NeoRL dataset trajectories & -282,885 $\pm$ 1,197 & -- & -- \\
BC & NeoRL simulator & -288,411 $\pm$ 267 & 0 & -- \\
M100: GRU + CEM & NeoRL simulator & -274,181 $\pm$ 396 & +14,229 & 100\% \\
M1000: Transformer-2L + MPPI & NeoRL simulator & -220,129 $\pm$ 1,699 & +68,282 & 100\% \\
M10000: Transformer-2L + MPPI & NeoRL simulator & -223,876 $\pm$ 1,033 & +64,535 & 100\% \\
\bottomrule
\end{tabular}}
\end{table}


最终M100、M1000和M10000系统相对BC的平均提升分别为14,229、68,282和64,535，win rate均为100\%。这直接回答了Task 3：在当前NeoRL Industrial Benchmark设置中，经过严格模型选择和策略冻结的World-Model-based optimization能够提高长期cumulative reward。
